\section*{Chapitre \ref{ch:b2:global-change} --- Changement global et biosphère}

\begin{solution}{exo:b2:global-change:1}
$\Delta T = \lambda C_{\text{cum}}$ avec $\lambda = \qty{1.65}{\celsius}$
par \qty{1000}{GtC}~: \qty{1.2}{\celsius}, \qty{1.65}{\celsius},
\qty{3.3}{\celsius}.
\end{solution}

\begin{solution}{exo:b2:global-change:2}
Degrés-jours~: la somme, jour après jour, de l'excès de la température
moyenne journalière sur une base, le \omterm{thm:b2:global-change:phenology}{temps thermique} qui gouverne le
développement des plantes et des ectothermes. Décalage~: lorsque deux
espèces en interaction règlent leur calendrier sur des signaux différents
et que le réchauffement déplace l'un et non l'autre --- le gobemouche,
réglé sur la durée du jour en Afrique, arrive à son ancienne date et
trouve le pic de chenilles, réglé sur la température, déjà passé depuis
deux semaines.
\end{solution}

\begin{solution}{exo:b2:global-change:3}
Environ \qty{150}{m} vers le haut et \qty{150}{km} vers le pôle par
degré. Une espèce de sommet n'a pas de terrain plus élevé~: son isotherme
passe au-dessus de la montagne et son habitat disparaît, tandis que la
surface de chaque bande diminue avec l'altitude au cours de la montée.
\end{solution}

\begin{solution}{exo:b2:global-change:4}
Le \ensuremath{\mathrm{CO_2}} dissous forme de l'acide carbonique, qui
libère des ions hydrogène, et le pH baisse proportionnellement au
logarithme de la pression de \ensuremath{\mathrm{CO_2}}. Les ions hydrogène
supplémentaires transforment le carbonate en bicarbonate, si bien que la
concentration en carbonate et l'\omterm{thm:b2:global-change:acid}{état de saturation} du carbonate de
calcium diminuent. Les organismes qui construisent de l'aragonite ---
coraux, ptéropodes, de nombreuses larves --- ou de la calcite ---
coccolithophoridés, foraminifères, mollusques --- doivent dépenser
davantage pour construire leurs coquilles et, sous la saturation, les
perdent.
\end{solution}

\begin{solution}{exo:b2:global-change:5}
\qty{1.5}{\celsius}~: $1.5/1.65\times 1000 = \qty{909}{GtC}$ au total,
\qty{209}{GtC} restantes, 21 ans. \qty{2}{\celsius}~:
\qty{1212}{GtC}, \qty{512}{GtC} restantes, 51 ans.
\end{solution}

\begin{solution}{exo:b2:global-change:6}
$\sqrt{2\times 150/0.12} = 50$ jours~; avancée $1.5/0.12 = 12.5$
jours.
\end{solution}

\begin{solution}{exo:b2:global-change:7}
\qty{3}{\celsius}~: les isothermes montent de $3000/6.5 = \qty{460}{m}$~;
la bande passe à \qtyrange{1260}{1860}{m}, mais la montagne s'arrête à
\qty{1800}{m}~: les \qty{60}{m} supérieurs de la bande sont perdus et le
reste se situe là où la surface vaut $2^{-460/300} = 0.35$ de ce qu'elle
était --- environ deux tiers de la surface perdus. \qty{5}{\celsius}~:
\qty{770}{m}~; la bande serait à \qtyrange{1570}{2170}{m}~; seule la
tranche \qtyrange{1570}{1800}{m} existe, un tiers de la largeur de la
bande, à $2^{-770/300} = 0.17$ de la surface --- moins d'un dixième de la
surface d'origine~: l'espèce a pratiquement disparu.
\end{solution}

\begin{solution}{exo:b2:global-change:8}
560~: pH $7.93$, ions hydrogène $\times 1.9$~; 800~: $7.79$, $\times
2.6$~; 1000~: $7.70$, $\times 3.1$.
\end{solution}

\begin{solution}{exo:b2:global-change:9}
$1 - (1 - h)^{0.25}$~: \qty{26}{\%}, \qty{53}{\%}, \qty{82}{\%}. La
perte immédiate est plus faible parce que le fragment restant abrite
encore des populations de la plupart des espèces~; elles sont trop
petites pour persister et s'éteignent au fil des décennies --- la \omterm{thm:b2:global-change:sar}{dette
d'extinction}.
\end{solution}

\begin{solution}{exo:b2:global-change:10}
$c = 2\sqrt{0.7\times 20} = \qty{7.5}{km/yr}$~; \qty{1000}{km} en
environ 130 ans. Diviser $r$ ou $D$ par deux divise $c$ par
$\sqrt 2$, à \qty{5.3}{km/yr}~; ni l'un ni l'autre ne la divise par deux
--- pour diviser la vitesse par deux, il faut diviser le produit $rD$ par
quatre.
\end{solution}

\begin{solution}{exo:b2:global-change:11}
Taux de fond~: $8\times 10^{6}\times 10^{-6} = 8$ par an, 800 par
siècle~; à mille fois ce taux, 8000 par an, \num{800000} par siècle. Dix
mille extinctions documentées en un siècle représentent environ dix fois
le taux de fond, et non mille --- mais la documentation ne couvre que les
quelques pour cent d'espèces décrites et suivies (oiseaux, mammifères,
quelques plantes), parmi lesquelles le taux vaut effectivement cent fois
le taux de fond ou plus~; pour la majorité non décrite, la plupart des
extinctions ne sont jamais observées, et l'estimation aire--espèces est
le seul moyen d'approche. La vérité se situe entre le compté et l'estimé,
d'où l'étendue de la fourchette citée.
\end{solution}

\begin{solution}{exo:b2:global-change:12}
La température est proportionnelle aux \omterm{prop:b2:global-change:tcre}{émissions cumulées}~; lorsque les
émissions nettes sont nulles, le total cumulé cesse de croître, et le
réchauffement aussi. Dans le \omterm{thm:b2:biogeochemical-cycles:box}{modèle à un compartiment}, avec des émissions
nulles, l'excédent atmosphérique commence à diminuer à mesure que les
puits continuent d'agir, ce qui refroidirait --- mais l'océan, qui se
réchauffe encore vers l'équilibre avec le \ensuremath{\mathrm{CO_2}} accru,
continue de réchauffer la surface, et les deux effets se compensent
presque~: la température reste où elle est pendant des siècles. Ce qui
manque à l'argument~: les autres gaz à \omterm{prop:b2:global-change:tcre}{effet de serre} (la courte durée de
vie du méthane fait que son réchauffement diminue vite une fois ses
émissions arrêtées), les aérosols qui masquent aujourd'hui une partie du
réchauffement et disparaissent quelques semaines après les émissions qui
les produisent, et les rétroactions --- carbone du pergélisol,
dépérissement des forêts --- qui pourraient ajouter leurs propres
émissions.
\end{solution}

\begin{solution}{pb:b2:global-change:1}
\textbf{1.} A~: $700 + 80\times 10 = \qty{1500}{GtC}$. B~: $700 +
\tfrac12\times 50\times 10 = \qty{950}{GtC}$.
\textbf{2.} A~: $1.65\times 1.5 = \qty{2.5}{\celsius}$. B~:
$1.65\times 0.95 = \qty{1.6}{\celsius}$.
\textbf{3.} 2050. A~: \qty{1000}{GtC}, \qty{1.65}{\celsius}. B~:
émissions $10\int_0^{30}(1 - t/50)\,\mathrm{d}t = 10(30 - 9) =
\qty{210}{GtC}$, total 910, \qty{1.5}{\celsius}.
\textbf{4.} $\lambda\times 100 = \qty{0.165}{\celsius}$ par décennie,
un peu moins que les $0.2$ observés (qui incluent les autres gaz et la
disparition du masquage par les aérosols).
\textbf{5.} Oui~: les \omterm{prop:b2:global-change:tcre}{émissions cumulées} cessent de croître en 2070, et le
réchauffement aussi, qui reste à \qty{1.6}{\celsius} --- la
proportionnalité signifie que la température est fixée par le total, non
par le rythme.
\textbf{6.} $2/1.65\times 1000 = \qty{1212}{GtC}$~; la trajectoire A les
atteint au bout de $512/10 = 51$ ans, en 2071.
\textbf{7.} $\sqrt{2\times 200/0.1} = 63$ jours.
\textbf{8.} Réchauffement par rapport à 2020~: A $2.5 - 1.2 =
\qty{1.3}{\celsius}$, B \qty{0.4}{\celsius}. Avancée $\Delta T/b$~:
A 13 jours, B 4 jours.
\textbf{9.} Écart~: A $20 + 13 = 33$ jours~; B 24 jours.
\textbf{10.} Chenilles disponibles du 15\textsuperscript{e} au
25\textsuperscript{e} jour après la feuillaison. B~: l'oiseau arrive au
24\textsuperscript{e} jour, dans la fenêtre, de justesse. A~: au
33\textsuperscript{e} jour, huit jours après la dernière chenille.
\textbf{11.} L'oiseau avance de $3\times 1.3 = 4$ jours~: écart $33 - 4 =
29$ jours --- toujours hors de la fenêtre.
\textbf{12.} L'oiseau peut vivre à la nouvelle température~; ce qu'il ne
peut pas faire, c'est nourrir ses oisillons une fois la nourriture
disparue, parce que son signal et celui de sa proie réagissent
différemment au réchauffement.
\textbf{13.} A~: $1.3/6.5\times 1000 = \qty{200}{m}$ vers le haut et
$1.3/0.65\times 100 = \qty{200}{km}$ vers le pôle. B~: \qty{57}{m} et
\qty{57}{km}.
\textbf{14.} En 80 ans, l'arbre se déplace de \qty{8}{km}~: il ne peut
pas suivre \qty{200}{km} (A) et reste aussi en deçà de \qty{57}{km} (B)
--- les arbres sont en retard, et leurs aires resteront en déséquilibre
avec le climat pendant des siècles sur l'une comme sur l'autre
trajectoire.
\textbf{15.} A~: la bande serait à \qtyrange{1700}{2000}{m}~: elle atteint
tout juste le sommet, sans nulle part où aller plus haut. La montagne se
rétrécit avec l'altitude, si bien que la même tranche ne couvre plus que
$2^{-200/300} = 0.63$ fois sa surface d'origine --- et tout
réchauffement supplémentaire n'a plus de terrain où la déplacer.
\textbf{16.} B~: \qty{57}{m} vers le haut, facteur de surface
$2^{-57/300} = 0.88$~: une perte de \qty{12}{\%}.
\textbf{17.} Huit décennies à \qty{30}{km}~: \qty{240}{km}, plus que les
\qty{200}{km} nécessaires sur A~: le poisson suit (et quitte sa pêcherie
d'origine).
\textbf{18.} À la limite inférieure, l'espèce rencontre des compétiteurs,
des pathogènes et des prédateurs montant d'en dessous auxquels elle
n'était pas confrontée auparavant, et sa propre physiologie, adaptée à
l'ancienne température, est surpassée par la leur~: le recul est
compétitif autant que physiologique.
\textbf{19.} A~: $\mathrm{pH} = 8.2 - 0.9\log_{10}(600/280) = 7.90$~;
ions hydrogène $10^{0.30} = 2.0$ fois ceux de 1750. B~: $8.04$~; $1.44$
fois.
\textbf{20.} $\Omega = 4/(\text{rapport})$~: A $2.0$, B $2.8$. À 2, les
coraux calcifient encore, mais lentement, leurs squelettes sont plus
fragiles et leurs récifs s'érodent plus vite qu'ils ne croissent~; à
$2.8$, ils construisent, même si la chaleur de la trajectoire B les fait
encore blanchir.
\textbf{21.} A~: $400\times 0.4^{0.25} = 318$, 82 perdues. B~:
$400\times 0.8^{0.25} = 378$, 22 perdues.
\textbf{22.} A~: $318\times(1 - 0.05\times 1.3) = 298$. B~:
$378\times(1 - 0.05\times 0.4) = 371$.
\textbf{23.} A~: $1 - 298/400 = \qty{26}{\%}$. B~: \qty{7}{\%}.
\textbf{24.} Sur les deux, le défrichement~: 82 espèces contre 20 sur A,
22 contre 7 sur B --- le terme de réchauffement du modèle est modeste~; en
réalité, les deux interagissent, car une forêt fragmentée ne peut pas
déplacer son aire.
\textbf{25.} Trajectoire A~: \qty{2.5}{\celsius}, pH $7.90$, un quart des
espèces d'arbres perdues. Trajectoire B~: \qty{1.6}{\celsius}, pH $8.04$,
\qty{7}{\%} perdues.
\end{solution}
